\section{Background and Related Work}\label{sec:b-related}

This section reviews the infrastructure the approach builds on, the analyses that already read quantum source code,
the tools that work on circuits and the dependence models proposed for slicing, and ends with the gaps that
Table~\ref{tab:b-related} summarizes.

\subsection{srcML and srcSlice}
srcML represents source code as XML in which the syntax of the language is marked up while the original text, comments
and layout included, is kept, so the XML converts back to the exact source~\cite{collard2011lightweight,
collard2013srcml}. Because srcML parses without compiling, it accepts fragments and programs whose libraries are
missing, which suits code collected from repositories. srcSlice computes slices from srcML for every variable of a
system at once~\cite{alomari2014srcslice}. Instead of building a program dependence graph~\cite{ferrante1987pdg,
horwitz1990interprocedural}, it keeps a slice profile per variable, the lines that define and use it, the variables it
flows into and the calls that receive it, and slices from these profiles, which let it process the Linux kernel.
Neither tool knows what a quantum program means. To srcML a Qiskit program is Python, and to srcSlice its circuit is a
single variable that every gate call uses, so the slice of any qubit is every line that touches the circuit.

\subsection{Analyses of quantum source code}
Behler et al.\ specified a srcML markup for OpenQASM~3~\cite{behler2024supporting} and built QStatic on it, which
derives the order in which a program applies operations to each qubit, unrolling loops and expanding gate definitions,
and uses that order to find iteration patterns and suggest refactorings~\cite{behler2025static}. They demonstrate the
markup on the 16 example files of OpenQASM~3, leave \texttt{while} loops for later, and name a slicer that finds the
lines affecting a qubit as future work. Kaul et al.\ translate Qiskit and OpenQASM programs into a code property graph,
the QCPG, and query it for bugs and code smells~\cite{kaul2023uniform}. LintQ runs CodeQL queries over Qiskit programs,
unrolling loops for up to ten iterations to relate their operations, and was applied to 7,568 programs from
GitHub~\cite{paltenghi2024lintq}. QChecker checks Qiskit programs against bug patterns over their syntax trees~\cite{zhao2023qchecker}, Q-PAC is
another pattern-based detector~\cite{nayak2023qpac}, Chen et al.\ define smells specific to quantum
code~\cite{chen2023smelly},
and Xia and Zhao analyze entanglement statically~\cite{xia2025entanglement}. These analyses look for known patterns and
properties. None returns the statements that determine a qubit.

\subsection{Tools that work on circuits}
Once a program has run and built its circuit, Qiskit can list the operations that precede a qubit's last operation as
the ancestors in its DAG representation, and its \texttt{LightCone} pass keeps the gates that can affect given qubits or
an observable and drops those that commute with it~\cite{qiskit2025}. Qiskit's slicing utilities split a circuit into consecutive pieces at barriers, by depth or by gate
type~\cite{qiskitaddonutils}, and the Cirquo tool of Metwalli and Van Meter splits a
circuit at barriers and maps its gates back to source lines through tracebacks taken while the program
runs~\cite{metwalli2022debugging, metwalli2024testing}. Working on the circuit has two costs for a developer. The
program must run, with its Qiskit version, its inputs and sometimes access to hardware, and the circuit no longer says
which line and which loop iteration added each gate.

\subsection{Dependence models for quantum programs}
Yu and Zhao describe a dependence graph for quantum programs whose edges add four quantum kinds to data and control
dependence, and outline slicing on it without implementing it~\cite{yu2025qpdg}. Compilers order and place gates by their backward light cone, the gates whose effect can propagate to given qubits by
the end of a circuit~\cite{murali2019triq, amy2020staq}. The approach slices by this light cone, followed further from each
measurement to the operations its bit conditions. Two roadmaps point to the same need, one describing dependence graphs
as an aid to understanding quantum programs~\cite{murillo2025roadmap}, the other finding that dependence analysis of
circuits has hardly been studied~\cite{ramalho2025road}, and developers report that their debugging methods, printing, drawing circuits and reading code, do not
scale~\cite{zappin2026challenges}.

\subsection{Gaps}
Table~\ref{tab:b-related} compares these approaches with \tool{}. Three gaps remain. No slicer reads Qiskit, the
language most quantum programs are written in, and the slicer named for OpenQASM~3 is future work. The tools that
slice or prune circuits need the program to run and lose the link to its lines and loop iterations. The analyzers that
read Qiskit source look for patterns, and none has been evaluated as a slicer on real programs.

\begin{table*}[t]
  \centering
  \caption{Approaches that analyze quantum programs or circuits, against what a source-level quantum slicer needs}
  \label{tab:b-related}
  \footnotesize\setlength{\tabcolsep}{3.4pt}\renewcommand{\arraystretch}{1.18}
  \newcommand{\sep}{\arrayrulecolor{black!25}\hline\arrayrulecolor{black}}
  \begin{tabularx}{\textwidth}{@{}>{\raggedright\arraybackslash}p{3.0cm} *{6}{>{\centering\arraybackslash}p{1.2cm}} >{\raggedright\arraybackslash}X@{}}
    \toprule
    \textbf{Approach} & \textbf{Reads Qiskit} & \textbf{Without running} & \textbf{Slice per qubit} & \textbf{Measured bits} &
    \textbf{Loops} & \textbf{Source lines} & \textbf{Evaluated on} \\
    \midrule
    srcSlice~\cite{alomari2014srcslice}                                  & \none & \full & \none & \na   & \full & \full & Linux kernel \\ \sep
    Behler et al., QStatic~\cite{behler2024supporting, behler2025static} & \none & \full & \none & \half & \half & \full & 16 OpenQASM~3 example files \\ \sep
    QPDG, Yu and Zhao~\cite{yu2025qpdg}                                  & \full & \full & \full & \full & ?     & \full & 5 listings, not implemented \\ \sep
    QCPG, Kaul et al.~\cite{kaul2023uniform}                             & \full & \full & \none & \full & \full & \full & 3 listings, 1 Qiskit/OpenQASM circuit \\ \sep
    LintQ~\cite{paltenghi2024lintq}                                      & \full & \full & \none & \full & \half & \full & 7,568 GitHub programs, 216 problems found by inspection \\ \sep
    QChecker~\cite{zhao2023qchecker}                                     & \full & \full & \none & \na   & \na   & \full & 42 Bugs4Q bugs \\ \sep
    Qiskit DAG, \texttt{LightCone}~\cite{qiskit2025}                     & \none & \half & \full & \none & \none & \none & \na \\ \sep
    Qiskit circuit slicing~\cite{qiskitaddonutils}                       & \none & \half & \none & \na   & \na   & \none & \na \\ \sep
    Cirquo, Metwalli and Van Meter~\cite{metwalli2022debugging, metwalli2024testing} & \full & \none & \none & \na & \na & \half & 1 Grover case study, 1 user survey \\
    \midrule
    \rowcolor{blue!7}\textbf{\tool{}} & \full & \full & \full & \full & \full & \full & \numSizeLintqFiles{} GitHub, \numSizeBugsFiles{} Bugs4Q, \numSizeTextFiles{} textbook and \numSizeQheFiles{} Qiskit HumanEval programs; \numScaleQbFiles{} QASMBench circuits \\
    \bottomrule
  \end{tabularx}\\[3pt]
  \begin{minipage}{\textwidth}\footnotesize
  \full{} has the property, \half{} has it in part, \none{} lacks it, ? means the authors do not specify it, and \na{}
  means it does not apply. \emph{Reads Qiskit}, the approach takes Qiskit programs as input. \emph{Without running}, it
  analyzes the program without executing it; circuit tools need the program run to obtain the circuit, unless the
  circuit is given in OpenQASM. \emph{Slice per qubit}, it returns the operations or lines that can affect a chosen
  qubit. \emph{Measured bits}, it follows a measured bit to the operations the bit conditions. \emph{Loops}, it follows
  loops of any length. \emph{Source lines}, it reports results on the program's lines. \emph{Evaluated on}, the data of
  the original evaluation.
  \end{minipage}
\end{table*}
